% Czech and fonts
\documentclass[11pt, a4paper]{article}
\usepackage[utf8]{inputenc} 
\usepackage[T1]{fontenc} 
\usepackage[czech]{babel}
\usepackage{lmodern}

% Margins and colors
\usepackage[top=2.5cm, bottom=2.5cm, left=3cm, right=2.5cm]{geometry}
\usepackage{graphicx} 
\usepackage{xcolor} 

% Code listings
\usepackage{listings} 
\lstset{
    backgroundcolor=\color{white},
    basicstyle=\ttfamily\small,
    breaklines=true,
    frame=none,
    xleftmargin=1cm,
    xrightmargin=0.5cm,
    aboveskip=4pt,
    belowskip=4pt,
}

% Hyperlinks 
\usepackage[hidelinks]{hyperref}

% Header and footer
\usepackage{fancyhdr} 
\pagestyle{fancy} 
\fancyhf{} 
\renewcommand{\headrulewidth}{0.4pt} 
\renewcommand{\footrulewidth}{0pt} 
\fancyhead[C]{\small\itshape Úloha č. 1 \quad Matematická kartografie} 
\fancyfoot[C]{\thepage}

% Spacing
\usepackage{setspace} 
\setstretch{1.5}

\usepackage{amsmath}
\usepackage{booktabs}
\usepackage{tabularx}

\begin{document}

% Titulní strana
\begin{titlepage}
    \centering

    {\large Přírodovědecká fakulta}\\[0.3cm]
    {\large Univerzita Karlova}\\[3cm]

    \includegraphics[width=0.425\textwidth]{imageUK.png}\\[3cm]

    {\LARGE\bfseries Výpočet souřadnic bodů v~rovině Křovákova zobrazení}\\[0.5cm]
    {\large Matematická kartografie}\\[4cm]

    {\large Kristýna Antošová, Vít Kadlec}\\[0.5cm]
    {\normalsize 1.N-GKDPZ}\\[0.3cm]
    {\normalsize Praha 2026}

    \vfill
\end{titlepage}

% Zadání
\section{Zadání}

\begin{figure}[h!]
\centering
\includegraphics[width=1\textwidth]{Zadani.png}
\label{fig:Zadani}
\end{figure}

\begin{figure}[h!]
\centering
\includegraphics[width=1\textwidth]{Hodnocení.png}
\label{fig:Hodnoceni}
\end{figure}
\noindent V rámci zadání byla řešena pouze základní úloha  \textit{Výpočet souřadnic bodu v rovině Křovákova zobrazení + zkreslení}. Bonusové úlohy řešeny nebyly.

\newpage

% Popis a rozbor problému
\section{Popis a~rozbor problému}

Transformace bodu $P' = [\varphi', \lambda']$ ze~systému zeměpisných souřadnic elipsoidu WGS-84 na bod $P = [Y, X]$ v~rovině Křovákova zobrazení (S-JTSK) probíhá v~následujících fázích:

\begin{enumerate}
    \item Převod zeměpisných souřadnic $[\varphi, \lambda]$ na geocentrické prostorové souřadnice $[X, Y, Z]$.
    \item Prostorová podobnostní (Helmertova) transformace mezi elipsoidy WGS-84 a~Bessel.
    \item Převod geocentrických souřadnic zpět na zeměpisné souřadnice Besselova elipsoidu.
    \item Gaussovo konformní zobrazení na~Gaussově kouli.
    \item Transformace do~šikmé polohy kartografického pólu.
    \item Lambertovo konformní kuželové zobrazení do~roviny.
\end{enumerate}

\subsection{Převod zeměpisných souřadnic na geocentrické}

Prostorové geocentrické souřadnice $(X, Y, Z)$ bodu na povrchu referenčního elipsoidu daného parametry $(a, b)$ lze určit ze~vztahu
\begin{equation}
    \mathbf{P} = \begin{pmatrix} X \\ Y \\ Z \end{pmatrix}
    = N \begin{pmatrix} \cos\varphi\cos\lambda \\ \cos\varphi\sin\lambda \\ (1-e^2)\sin\varphi \end{pmatrix},
\end{equation}
kde $N$ je příčný poloměr křivosti
\begin{equation}
    N = \frac{a}{W}, \qquad W = \sqrt{1 - e^2 \sin^2\varphi},
\end{equation}
a~$e^2$ je první excentricita elipsoidu
\begin{equation}
    e^2 = \frac{a^2 - b^2}{a^2}.
\end{equation}

\subsection{Helmertova prostorová podobnostní transformace}

Transformace mezi geocentrickými souřadnicemi WGS-84 a~Besselova elipsoidu je dána sedmiprvkovou Helmertovou transformací
\begin{equation}
    \begin{pmatrix} X \\ Y \\ Z \end{pmatrix}_{\mathrm{Bes}}
    = m \begin{pmatrix} 1 & \omega_z & -\omega_y \\ -\omega_z & 1 & \omega_x \\ \omega_y & -\omega_x & 1 \end{pmatrix}
    \begin{pmatrix} X \\ Y \\ Z \end{pmatrix}_{\mathrm{WGS}}
    + \begin{pmatrix} \Delta X \\ \Delta Y \\ \Delta Z \end{pmatrix},
\end{equation}
kde matice rotace využívá linearizaci $\sin\omega \approx \omega$, $\cos\omega \approx 1$ pro malé hodnoty rotací. Parametry transformace jsou:

\begin{table}[h!]
\centering
\begin{tabular}{l r}
\toprule
\textbf{Parametr} & \textbf{Hodnota} \\
\midrule
$\omega_x$ & $4{,}9984''$ \\
$\omega_y$ & $1{,}5867''$ \\
$\omega_z$ & $5{,}2611''$ \\
$m-1$ & $-3{,}5623 \times 10^{-6}$ \\
$\Delta X$ & $-570{,}8285$ m \\
$\Delta Y$ & $-85{,}6769$ m \\
$\Delta Z$ & $-462{,}8420$ m \\
\bottomrule
\end{tabular}
\end{table}

\subsection{Převod geocentrických souřadnic na zeměpisné}

Zeměpisná délka se určí ze~vztahu
\begin{equation}
    \tan\lambda = \frac{Y}{X}.
\end{equation}
Zeměpisná šířka se získá z~aproximace
\begin{equation}
    \tan\varphi = \frac{Z}{(1-e^2)\sqrt{X^2 + Y^2}}.
\end{equation}

\subsection{Gaussovo konformní zobrazení Besselova elipsoidu na kouli}

Zeměpisná délka se převede na délku od Ferra: $\lambda_{\mathrm{Ferro}} = \lambda_{\mathrm{Bes}} + 17^{\circ}40'$. Následně se bod zobrazí na Gaussovu kouli o~poloměru $R$ pomocí sférických souřadnic $(u, v)$:
\begin{equation}
    \tan\left(\frac{u}{2} + \frac{\pi}{4}\right)
    = \frac{1}{k} \left[
    \tan\left(\frac{\varphi}{2} + \frac{\pi}{4}\right)
    \cdot \left(\frac{1 - e\sin\varphi}{1 + e\sin\varphi}\right)^{e/2}
    \right]^\alpha, \qquad
    v = \alpha \cdot \lambda_{\mathrm{Ferro}},
\end{equation}
kde $\alpha$, $k$, $R$ jsou konstanty odvozené z~parametrů Besselova elipsoidu a~základní rovnoběžky $\varphi_0 = 49{,}5^{\circ}$.

\subsection{Transformace do~šikmé polohy}

Bod $(u, v)$ na Gaussově kouli se transformuje do~šikmé polohy $(s, d)$ vzhledem ke~kartografickému pólu $K = [u_k, v_k]$. Ze~sférické trigonometrie plyne:
\begin{equation}
    \sin s = \sin u \cdot \sin u_k + \cos u \cdot \cos u_k \cdot \cos(v_k - v),
\end{equation}
\begin{equation}
    \tan d = \frac{\sin(v_k - v) \cdot \cos u}{-\sin u \cdot \cos u_k + \cos u \cdot \cos(v_k - v) \cdot \sin u_k}.
\end{equation}
Kartografický pól Křovákova zobrazení má souřadnice $u_k = 59^{\circ}42'42{,}6969''$, $v_k = 42^{\circ}31'31{,}41725''$.

\subsection{Lambertovo konformní kuželové zobrazení}

V~závěrečném kroku se šikmé souřadnice zobrazí do~roviny Křovákova zobrazení:
\begin{equation}
    \rho = \rho_0 \cdot \left(\frac{\tan(s_0/2 + \pi/4)}{\tan(s/2 + \pi/4)}\right)^c, \qquad
    \varepsilon = c \cdot d,
\end{equation}
kde $s_0 = 78{,}5^{\circ}$, $c = \sin s_0$, $\rho_0 = 0{,}9999 \cdot R / \tan s_0$. Výsledné rovinné souřadnice jsou:
\begin{equation}
    Y = \rho \sin\varepsilon, \qquad X = \rho \cos\varepsilon.
\end{equation}

\subsection{Délkové zkreslení}

Délkové zkreslení $m$ v~daném bodě je dáno vztahem
\begin{equation}
    m = \frac{c \cdot \rho}{R \cdot \cos s},
\end{equation}
přičemž hodnota $m-1$ vyjadřuje relativní odchylku od~nezkreslené délky.

\newpage

% Výsledky
\section{Výsledky}

Výsledky transformace obou bodů shrnuje tabulka~\ref{tab:vysledky}. Vstupem jsou souřadnice změřené pomocí integrovaného GPS modulu v mobilním telefonu na dvou trigonometrických bodech v~systému WGS-84. Výstupem jsou rovinné souřadnice v~S-JTSK a~délkové zkreslení.

\begin{table}[h!]
\centering
\caption{Výsledky transformace WGS-84 $\rightarrow$ S-JTSK}
\label{tab:vysledky}
\begin{tabularx}{\textwidth}{X r r}
\toprule
\textbf & \textbf{Bod $P_1$} & \textbf{Bod $P_2$} \\
\midrule
\multicolumn{3}{l}{\textit{Vstup (WGS-84)}} \\
$\varphi$ & $50{,}064270^{\circ}$ & $50{,}050572^{\circ}$ \\
$\lambda$ & $14{,}419012^{\circ}$ & $14{,}384470^{\circ}$ \\
\midrule
\multicolumn{3}{l}{\textit{Výstup (S-JTSK)}} \\
$Y$ [m] & $743\,320{,}365$ & $745\,978{,}602$ \\
$X$ [m] & $1\,045\,546{,}844$ & $1\,046\,718{,}237$ \\
\midrule
\multicolumn{3}{l}{\textit{Délkové zkreslení}} \\
$m - 1$ & $-9{,}715 \times 10^{-5}$ & $-9{,}801 \times 10^{-5}$ \\
\bottomrule
\end{tabularx}
\end{table}

\noindent Záporné hodnoty délkového zkreslení $m-1$ znamenají, že délky v~rovině zobrazení jsou nepatrně kratší než odpovídající délky na~elipsoidu, což odpovídá vlastnostem Křovákova zobrazení v~oblasti Prahy.

\subsection{Vizualizace}

Obrazy bodů $P_1$ a~$P_2$ byly vizualizovány s~podkladem mapy OpenStreetMap. Na~každém výkresu je vyznačen vstupní bod (bílý kroužek) a~jeho transformovaný obraz v~S-JTSK (modrý trojúhelník). Oba symboly leží na~stejném fyzickém místě, protože transformace zachovává polohu bodu.

\begin{figure}[h!]
\centering
\includegraphics[width=0.7\textwidth]{P1_map.png}
\caption{Vizualizace bodu $P_1$ -- WGS-84 a~S-JTSK.}
\label{fig:p1}
\end{figure}

\begin{figure}[h!]
\centering
\includegraphics[width=0.7\textwidth]{P2_map.png}
\caption{Vizualizace bodu $P_2$ -- WGS-84 a~S-JTSK.}
\label{fig:p2}
\end{figure}

\clearpage

% Skripty
\section{Skripty}

Výpočty a~vizualizace byly realizovány v~jazyce Python.

\begin{itemize}
    \item \texttt{uvtosd.py} --- transformace bodu $[u, v]$ do~šikmé polohy zobrazení $[s, d]$ vzhledem ke~kartografickému pólu $K = [u_k, v_k]$, implementace rovnic~(8) a~(9).
    \item \texttt{u1.py} --- hlavní skript obsahující:
    \begin{itemize}
        \item funkci \texttt{WGSToJTSK} realizující celý transformační řetězec WGS-84 $\rightarrow$ S-JTSK dle rovnic~(1)--(11) včetně výpočtu délkového zkreslení~(12),
        \item funkce \texttt{fetch\_osm\_tiles}, \texttt{\_deg2tile}, \texttt{\_tile2deg} pro stažení a~složení mapových dlaždic z~OpenStreetMap, funkce byly vytvořeny s pomocí modelu Claude Opus 4.6
        \item funkci \texttt{visualize} pro vykreslení bodu na~mapovém podkladu. Funkce byla vytvořena s pomocí modelu Claude Opus 4.6
    \end{itemize}
\end{itemize}

\newpage

% Zhodnocení a závěr
\section{Zhodnocení a~závěr}

Byly vypočteny souřadnice bodů $P_1$ a~$P_2$ v~rovině Křovákova zobrazení (S-JTSK) ze~zeměpisných souřadnic změřených GPS v~systému WGS-84. Transformační řetězec zahrnuje převod na~geocentrické souřadnice, sedmiprvkovou Helmertovu transformaci na~Besselův elipsoid, Gaussovo konformní zobrazení na~sféru, transformaci do~šikmé polohy a~Lambertovo konformní kuželové zobrazení.

Správnost výsledků byla ověřena porovnáním s~údaji z~databáze ČÚZK, přičemž odchylka nepřesahuje 1~m, což odpovídá přesnosti GPS měření. Délkové zkreslení $m-1$ dosahuje v~obou bodech záporných hodnot řádu $10^{-5}$, tedy délky v~rovině zobrazení jsou zkráceny přibližně o~0{,}01~\%.

\newpage
\renewcommand{\refname}{Zdroje} 
\begin{thebibliography}{9}

\bibitem{bayer_lcc} BAYER, T. (2026): \textit{Jednoduchá kuželová zobrazení. Křovákovo zobrazení}. Přednáška pro předmět Matematická kartografie, Katedra aplikované geoinformatiky a~kartografie, Přírodovědecká fakulta UK [cit.~11.5.2026].

\bibitem{bayer_navod} BAYER, T. (2026): \textit{Transformace zeměpisných souřadnic mezi elipsoidy} (stručný návod k~úloze~1). Výukový materiál pro předmět Matematická kartografie, Katedra aplikované geoinformatiky a~kartografie, Přírodovědecká fakulta UK [cit.~11.5.202].

\bibitem{osm} OpenStreetMap contributors (2026): \textit{OpenStreetMap}. \url{https://www.openstreetmap.org/} [cit.~11.5.2026]. Data dostupná pod licencí ODbL.

\end{thebibliography}
\end{document}